% c01a 原书 p.1–4（PDF p.008–011）
\subsection{前层与层}
\label{sec:1.1}

设 $X$ 是拓扑空间。$X$ 上集合的\emph{前层}\idx{qianceng}{前层} $\cF$ 由以下数据组成：
\begin{enumerate}[label=(\alph*),leftmargin=2.8em]
  \item 对 $X$ 的每个非空开子集 $U$，有一个集合 $\cF(U)$，其元素称为 $\cF$ 在 $U$ 上的\emph{截面}\idx{jiemian}{截面}；
  \item 对 $X$ 的非空开子集的每个包含 $V\subset U$，有一个映射 $\rho_{UV}\colon \cF(U)\to \cF(V)$，称为\emph{限制}\idx{xianzhi}{限制}。
\end{enumerate}
这些数据满足以下性质：
\begin{enumerate}[label=(\roman*),leftmargin=2.8em]
  \item $\rho_{UU}=\id_{\cF(U)}$。
  \item 若 $W\subset V\subset U$ 是 $X$ 的开子集，则 $\rho_{UW}=\rho_{VW}\rho_{UV}$。
\end{enumerate}

对任意截面 $s\in \cF(U)$ 与 $V\subset U$，常把 $\rho_{UV}(s)$ 记为 $s|_V$。为在记号中显式标出开集 $U$，$\cF(U)$ 的元素常记为 $(s,U)$。

$X$ 的开子集构成范畴，故对任意两个开子集 $U$ 与 $V$，有
\[
\Hom(V,U)=
\begin{cases}
\emptyset, & \text{若 } V\not\subset U,\\
\{V\to U\}, & \text{若 } V\subset U.
\end{cases}
\]
$X$ 上集合的前层 $\cF$ 就是从 $X$ 的开子集范畴到集合范畴的反变函子。

类似地，定义 $X$ 上交换群（resp.\ 环）的前层 $\cF$ 为从 $X$ 的开子集范畴到交换群（resp.\ 环）范畴的反变函子。我们约定：当 $U$ 为空集时，$\cF(U)=\{0\}$。

下面是前层的一些例子：
\begin{enumerate}[label=\arabic*.,leftmargin=2.2em,itemsep=0.2em]
  \item 设 $X$ 是拓扑空间，$A$ 是交换群。对 $X$ 的每个非空开子集 $U$，定义 $\cF(U)=A$；对非空开子集的每个包含 $V\subset U$，定义 $\rho_{UV}=\id_A$。则 $\cF$ 是交换群的前层，称为与 $A$ 相伴的\emph{常前层}\idx{changqianceng}{常前层}。
  \item 设 $X$ 是拓扑空间。对 $X$ 的每个非空开子集 $U$，定义 $C(U)$ 为 $U$ 上复值连续函数构成的环；对非空开子集的每个包含 $V\subset U$，定义 $\rho_{UV}\colon C(U)\to C(V)$ 为函数的限制。则 $C$ 是环的前层。
  \item 设 $\pi\colon X'\to X$ 是拓扑空间的连续映射。对 $X$ 的每个非空开子集 $U$，定义 $\cS(U)$ 为 $\pi$ 在 $U$ 上的连续截面构成的集合：
  \[
  \cS(U)=\{s\colon U\to \pi^{-1}(U)\mid \pi s=\id,\ \text{且 } s \text{ 连续}\},
  \]
  并对非空开子集的每个包含 $V\subset U$，定义 $\rho_{UV}\colon \cS(U)\to \cS(V)$ 为截面的限制。则 $\cS$ 是集合的前层。
\end{enumerate}

设 $\cF$ 是 $X$ 上集合（resp.\ 交换群，resp.\ 环）的前层。若它满足以下两个条件，则称 $\cF$ 为\emph{层}\idx{ceng}{层}：
\begin{enumerate}[label=(\roman*),leftmargin=2.8em]
  \item 设 $s,t\in \cF(U)$ 是两个截面。若存在 $U$ 的开覆盖 $\{U_i\}_{i\in I}$，使 $s|_{U_i}=t|_{U_i}$ 对所有 $i$ 成立，则 $s=t$。
  \item 设 $\{U_i\}_{i\in I}$ 是 $U$ 的开覆盖，$s_i\in \cF(U_i)$ 是截面，满足 $s_i|_{U_i\cap U_j}=s_j|_{U_i\cap U_j}$ 对所有 $i,j\in I$ 成立。则存在截面 $s\in \cF(U)$，使 $s|_{U_i}=s_i$ 对所有 $i\in I$ 成立。（由 (i)，这样的 $s$ 唯一。）
\end{enumerate}

注意，交换群的前层 $\cF$ 是层，当且仅当对任意开子集 $U$ 的任意开覆盖 $\{U_i\}_{i\in I}$，序列
\[
0\to \cF(U)\to \prod_{i\in I}\cF(U_i)\to \prod_{i,j\in I}\cF(U_i\cap U_j)
\]
正合，其中第二个箭头是
\[
\cF(U)\to \prod_{i\in I}\cF(U_i),\quad s\mapsto (s|_{U_i}),
\]
第三个箭头是
\[
\prod_{i\in I}\cF(U_i)\to \prod_{i,j\in I}\cF(U_i\cap U_j),\quad (s_i)\mapsto (s_j|_{U_i\cap U_j}-s_i|_{U_i\cap U_j}).
\]

上面最后两个前层的例子都是层。

\emph{有向集}\idx{youxiangji}{有向集}是指偏序集 $(I,\le)$，使得对任意 $i,j\in I$，存在 $k\in I$ 使 $i,j\le k$。集合的\emph{有向系统}\idx{youxiangxitong}{有向系统} $(A_i,\phi_{ij})_{i\in I}$ 由一族集合 $A_i$（$i\in I$）与映射 $\phi_{ij}\colon A_i\to A_j$（对 $i\le j$）组成，满足 $\phi_{ii}=\id_{A_i}$，且当 $i\le j\le k$ 时 $\phi_{jk}\phi_{ij}=\phi_{ik}$。对任意 $x_i\in A_i$ 与 $x_j\in A_j$，若存在 $k\ge i,j$ 使 $\phi_{ik}(x_i)=\phi_{jk}(x_j)$，则称 $x_i$ 与 $x_j$ \emph{等价}。这在不交并 $\coprod_{i\in I}A_i$ 上定义了等价关系。\emph{有向极限}\idx{youxiangjixian}{有向极限} $\dlim_i A_i$ 定义为 $(A_i,\phi_{ij})_{i\in I}$ 的等价类构成的集合。有典范映射 $A_i\to \dlim_i A_i$，把每个 $x\in A_i$ 送到与 $x$ 相应的等价类。
\begin{tnote}
本书把 direct limit 译作「有向极限」（中文文献亦作「有向极限」「归纳极限」），把 inverse limit 译作「逆向极限」；两者在本书中一律作此译法。[来源：本库 SY07-08 术语表, 工程文件/glossary.tsv]
\end{tnote}

设 $X$ 是拓扑空间，$P$ 是 $X$ 中的点。对 $P$ 的任意两个邻域 $U$ 与 $V$，若 $U\subset V$ 则称 $V\le U$。于是 $P$ 的邻域族按此序是有向集。对 $X$ 上任意前层 $\cF$，定义 $\cF$ 在 $P$ 处的\emph{茎}\idx{jing}{茎} $\cF_P$ 为
\[
\cF_P=\dlim_{P\in U}\cF(U),
\]
其中有向极限取遍 $P$ 的邻域族。$\cF_P$ 的元素可由 $\cF$ 在 $P$ 的某邻域上的截面代表。两个截面 $s\in \cF(U)$ 与 $t\in \cF(V)$ 定义 $\cF_P$ 中的同一元素，当且仅当存在 $P$ 的邻域 $W$，使 $W\subset U\cap V$ 且 $s|_W=t|_W$。对 $P$ 的每个邻域 $U$，有典范映射 $\cF(U)\to \cF_P$。截面 $s\in \cF(U)$ 在 $\cF_P$ 中的像称为 $s$ 在 $P$ 处的\emph{芽}\idx{ya}{芽}，记作 $s_P$。

设 $\cF$ 与 $\cG$ 是 $X$ 上交换群的前层。前层的\emph{态射}\idx{taishe}{态射} $\phi\colon \cF\to \cG$ 由一族交换群同态 $\phi(U)\colon \cF(U)\to \cG(U)$（$U$ 取遍开子集）组成，使得对开子集的每个包含 $V\subset U$，下图交换：
\begin{center}
\begin{tikzcd}
\cF(U) \arrow[r,"\phi(U)"] \arrow[d,"\rho_{UV}"'] & \cG(U) \arrow[d,"\rho_{UV}"] \\
\cF(V) \arrow[r,"\phi(V)"'] & \cG(V).
\end{tikzcd}
\end{center}

对任意点 $P\in X$，$\phi$ 诱导茎上的同态 $\phi_P\colon \cF_P\to \cG_P$。若把交换群的前层视为从 $X$ 的开子集范畴到交换群范畴的反变函子，则前层的态射恰是自然变换。类似地可定义集合或环的前层之间的态射。对任意前层 $\cF$，有恒等态射 $\id_{\cF}$。给定两个前层态射 $\phi\colon \cF\to \cG$ 与 $\psi\colon \cG\to \cH$，可按显然的方式定义它们的复合 $\psi\phi\colon \cF\to \cH$。于是得到前层范畴。前层态射 $\phi\colon \cF\to \cG$ 称为\emph{同构}\idx{tonggou}{同构}，若它有双边逆，即存在前层态射 $\psi\colon \cG\to \cF$，使 $\psi\phi=\id_{\cF}$ 且 $\phi\psi=\id_{\cG}$。这等价于说，$\phi(U)\colon \cF(U)\to \cG(U)$ 对每个开子集 $U$ 都是同构。我们把层之间的态射定义为它们作为前层的态射。于是得到层范畴，它是前层范畴的满子范畴。

\begin{proposition}\label{it:1.1.1}
设 $\phi\colon \cF\to \cG$ 是拓扑空间 $X$ 上层之间的态射。则 $\phi$ 是同构，当且仅当诱导的茎映射 $\phi_P\colon \cF_P\to \cG_P$ 对每个 $P\in X$ 都是同构。
\end{proposition}

\begin{proof}
「仅当」部分显然。下面证明「当」部分。设 $\phi_P\colon \cF_P\to \cG_P$ 对每个 $P\in X$ 都是双射。要证 $\phi(U)\colon \cF(U)\to \cG(U)$ 对 $X$ 的每个开子集 $U$ 都是双射。

设 $s,s'\in \cF(U)$ 是两个截面，使 $\phi(s)=\phi(s')$。则对每个 $P\in U$，有 $\phi_P(s_P)=\phi_P(s'_P)$。由 $\phi_P$ 是单射，得 $s_P=s'_P$。于是存在 $P$ 的邻域 $U_P\subset U$，使 $s|_{U_P}=s'|_{U_P}$。注意 $\{U_P\}_{P\in U}$ 是 $U$ 的开覆盖。因 $\cF$ 是层，必有 $s=s'$。所以 $\phi(U)$ 是单射。

设 $(t,U)$ 是 $\cG(U)$ 中的截面。对每个 $P\in U$，因 $\phi_P\colon \cF_P\to \cG_P$ 是满射，存在芽 $s_P\in \cF_P$ 使 $\phi_P(s_P)=t_P$。可设 $s_P$ 是某截面 $(s,U_P)\in \cF(U_P)$ 的芽，其中 $U_P$ 是 $P$ 的某邻域。注意 $\phi(s,U_P)$ 与 $(t,U)$ 在 $P$ 处有相同的芽。取 $U_P$ 足够小，可设 $U_P\subset U$ 且 $\phi(s,U_P)=(t,U)|_{U_P}$。于是对 $U$ 中任意两点 $P,Q$，有
\[
\phi(s,U_P)|_{U_P\cap U_Q}=(t,U)|_{U_P\cap U_Q}=\phi(s,U_Q)|_{U_P\cap U_Q}.
\]
由上面已证的 $\phi(U_P\cap U_Q)$ 是单射，必有 $(s,U_P)|_{U_P\cap U_Q}=(s,U_Q)|_{U_P\cap U_Q}$。注意 $\{U_P\}_{P\in U}$ 构成 $U$ 的开覆盖。因 $\cF$ 是层，存在截面 $(s,U)\in \cF(U)$，使 $(s,U)|_{U_P}=(s,U_P)$ 对每个 $P\in U$ 成立。有 $\phi(s,U)|_{U_P}=(t,U)|_{U_P}$。因 $\cG$ 是层，必有 $\phi(s,U)=(t,U)$。所以 $\phi(U)$ 是满射。
\end{proof}

在继续之前，我们引入范畴论中的一些概念。回忆：范畴 $\sC$ 由一个对象类 $\operatorname{ob}\sC$ 组成；对 $\operatorname{ob}\sC$ 中任意对象 $A$ 与 $B$，有一个态射集 $\Hom(A,B)$。态射可以复合：
\[
\Hom(A,B)\times\Hom(B,C)\to \Hom(A,C),\quad (f,g)\mapsto gf,
\]
且复合满足结合律。此外，对每个 $A\in \operatorname{ob}\sC$ 有恒等态射 $\id_A\in \Hom(A,A)$，使
\[
f\id_A=\id_B f=f
\]
对 $\Hom(A,B)$ 中每个态射 $f\colon A\to B$ 成立。

$\Hom(A,B)$ 中的态射 $f\colon A\to B$ 称为\emph{单态射}或\emph{单射}\idx{danshe}{单态射}，若对任意两个态射 $\alpha,\beta\colon C\to A$ 满足 $f\alpha=f\beta$，有 $\alpha=\beta$。满态射类似地由把箭头的方向反过来定义。更确切地说，$\Hom(B,A)$ 中的态射 $f\colon B\to A$ 称为\emph{满态射}或\emph{满射}\idx{manshe}{满态射}，若对任意两个态射 $\alpha,\beta\colon A\to C$ 满足 $\alpha f=\beta f$，有 $\alpha=\beta$。若态射既单又满，则称它为\emph{双射}\idx{shuangshen}{双射}。\emph{同构}是有双边逆的态射。任何同构都是双射。但双射态射未必是同构。
